% =====================================================================================================
\section{Protocol and Defence Evaluation}\label{sec:protocol}
% generated by paper/make_assets.py from paper/data -- do not edit
\begin{table}[t]
\caption{Reaction to each active attack (full suite, threshold policy, 3 runs; mean\,$\pm$\,s.d. in seconds). ``First rejection'' is when the receiver first discards a frame; ``mitigated'' is when the sender has raised protection or rekeyed.}
\label{tab:attacks}
\centering\small\setlength{\tabcolsep}{4pt}
\resizebox{\columnwidth}{!}{%
\begin{tabular}{lccc}
\toprule
Attack & Mitigated & First rejection & Mitigated after \\
\midrule
bit flips in the scan & 3/3 & 0.09\,$\pm$\,0.03 & 0.41\,$\pm$\,0.02 \\
mode downgrade in header & 3/3 & 0.06\,$\pm$\,0.02 & 0.44\,$\pm$\,0.02 \\
splice of two frames & 3/3 & 0.13\,$\pm$\,0.00 & 0.46\,$\pm$\,0.01 \\
forged frame injection & 3/3 & 0.14\,$\pm$\,0.06 & 0.44\,$\pm$\,0.03 \\
frame replay & 3/3 & 0.07\,$\pm$\,0.05 & 0.46\,$\pm$\,0.03 \\
frame delay & 3/3 & 4.13\,$\pm$\,0.04 & 1.46\,$\pm$\,0.03 \\
frame dropping & 3/3 & -- & 1.46\,$\pm$\,0.03 \\
junk flood & 3/3 & -- & 0.49\,$\pm$\,0.04 \\
feedback dropping & 3/3 & -- & 3.47\,$\pm$\,0.05 \\
forged feedback & 3/3 & -- & 0.46\,$\pm$\,0.04 \\
feedback replay & 3/3 & -- & 0.44\,$\pm$\,0.01 \\
\bottomrule
\end{tabular}}
\end{table}

\textit{Setup.} Sender, network channel with the attacker, and receiver run as separate processes over UDP, so the evaluation exercises the real wire format, fragmentation and timing (320$\times$240, 10 frames/s, datagrams of at most 1400 bytes). The full scenario applies \NumAttacks{} attacks for six seconds each, 25\,s apart: bit flips in the scan, a mode downgrade, splicing, injection of a forged frame, replay, delay, dropping, a junk flood, and dropping, forging and replaying feedback.

\textit{Results.} Over \SuiteRuns{} runs every attack was mitigated, on average \TtmMean{}\,$\pm$\,\TtmSd{}\,s after it began, with \SuiteFalse{} false alarms in calm windows (Table~\ref{tab:attacks}). With the PYNQ-Z2 as the sender, the same suite was again mitigated in full (\BdSuiteMit, column ``Board'') without a false alarm in \BdSuiteFalse{} calm windows, each attack within about one decision window. The evidence, not the harm, sets the reaction time: a bit flip fails the tag at the first frame, held-back frames produce no rejection while they are held, and silenced feedback blinds the sender until the dead-man rule fires after \FbDropTtm\,s. Integrity attacks are therefore caught at once; attacks by absence only by a timeout, whose length trades against false alarms on a slow network.

\textit{Benign impairments.} A second scenario applies only network problems: packet loss rising to 6\,\%, delay with jitter, and reordering. Over \BenignRuns{} runs the defence raised no false alarm (\BenignFalse{} windows), and none on the board (\BdBenignFalse{}). At \CaseSixPct\,\% packet loss the receiver lost \CaseSixFrameLoss\,\% of the frames, because a frame is lost when any of its fragments is; loss between 5 and 90\,\% is reported as a degraded network, not an attack, so jamming that lets more than a tenth of the frames through is indistinguishable from a bad network.

\textit{Fault injection inside the camera.} A variant of the two-lane board image places a fault injector between the encoder and the receive DMA, after the tag is computed, which flips one non-marker bit of the scan on a trigger. It stands for a fault or an implant anywhere on the path from the DMA through memory, the ARM and the network card. All \BdTamperShots{} triggered flips were detected and no other frame was rejected; \BdTamperOneFrame{} triggers cost exactly one frame and \BdTamperTwoFrames{} cost two, most likely because the trigger input is not debounced.
